\documentclass[12pt,a4paper]{article}

% -------------------------------------------------
% Sprache / Kodierung
% -------------------------------------------------
\usepackage[T1]{fontenc}
\usepackage[utf8]{inputenc}
\usepackage[ngerman]{babel}
\usepackage{lmodern}
\usepackage{microtype}

% -------------------------------------------------
% Mathematik
% -------------------------------------------------
\usepackage{amsmath,amssymb,amsthm,mathtools}

% -------------------------------------------------
% Layout / Grafik / Farben
% -------------------------------------------------
\usepackage[a4paper,margin=2.3cm]{geometry}
\usepackage{booktabs}
\usepackage{xcolor}
\usepackage[hidelinks]{hyperref}
\usepackage[most]{tcolorbox}
\usepackage{enumitem}

% -------------------------------------------------
% Lean-Code-Listings
% -------------------------------------------------
\usepackage{listings}
\usepackage{listingsutf8}

\definecolor{leanKeyword}{RGB}{0,70,170}
\definecolor{leanComment}{RGB}{40,120,40}
\definecolor{leanString}{RGB}{160,0,0}
\definecolor{leanBg}{RGB}{246,248,252}

\lstdefinelanguage{lean}{
  keywords={def,theorem,lemma,example,by,have,show,fun,let,where,import,
            open,namespace,section,end,variable,instance,class,structure,
            inductive,match,with,if,then,else,do,return,sorry,omega,simp,
            rfl,exact,apply,intro,intros,constructor,cases,rcases,
            obtain,use,ring,norm_num,decide,native_decide,linarith,
            Nat,Int,ZMod,Bool,Prop,Type,forall,exists},
  sensitive=true,
  morecomment=[l]{--},
  morecomment=[s]{/-}{-/},
  morestring=[b]",
  keywordstyle=\color{leanKeyword}\bfseries,
  commentstyle=\color{leanComment}\itshape,
  stringstyle=\color{leanString},
  basicstyle=\ttfamily\small,
  backgroundcolor=\color{leanBg},
  frame=single,
  framesep=4pt,
  xleftmargin=6pt,
  xrightmargin=6pt,
  breaklines=true,
  showstringspaces=false,
  tabsize=2,
  literate=
    {>=}{{$\geq$}}{1}
    {<=}{{$\leq$}}{1}
    {!=}{{$\neq$}}{1}
    {->}{{$\to$}}{1}
    {<->}{{$\leftrightarrow$}}{2},
}

\lstset{language=lean}

% -------------------------------------------------
% tcolorbox-Stile
% -------------------------------------------------
\definecolor{lückenrot}{RGB}{180,40,40}
\definecolor{bewiesengrün}{RGB}{30,100,50}
\definecolor{warnorange}{RGB}{200,100,10}
\definecolor{hintergrundblau}{RGB}{235,242,250}
\definecolor{hintergrundgelb}{RGB}{255,252,225}

\tcbset{
  lücke/.style={colback=red!8, colframe=lückenrot, fonttitle=\bfseries,
    title={Offene Lücke / Benötigt \texttt{sorry}}},
  kernidee/.style={colback=hintergrundblau, colframe=blue!60!black,
    fonttitle=\bfseries, title={Kernidee}},
  ergebnis/.style={colback=green!7, colframe=bewiesengrün,
    fonttitle=\bfseries, title={Beweisbar / In Mathlib verfügbar}},
  warnung/.style={colback=orange!10, colframe=warnorange,
    fonttitle=\bfseries, title={Achtung: Heuristik / Vermutung}},
  lean/.style={colback=leanBg!80, colframe=blue!40!black,
    fonttitle=\bfseries, title={Lean\,4-Code}},
}

% -------------------------------------------------
% Theorem-Umgebungen
% -------------------------------------------------
\newtheorem{definition}{Definition}[section]
\newtheorem{proposition}[definition]{Proposition}
\newtheorem{theorem}[definition]{Theorem}
\newtheorem{lemma}[definition]{Lemma}
\newtheorem{korollar}[definition]{Korollar}
\newtheorem{bemerkung}[definition]{Bemerkung}

% -------------------------------------------------
% Makros
% -------------------------------------------------
\newcommand{\vii}{\nu_2}
\newcommand{\U}{\mathcal{U}}
\newcommand{\N}{\mathbb{N}}
\newcommand{\Z}{\mathbb{Z}}
\newcommand{\Ztwo}{\mathbb{Z}_2}
\newcommand{\Syr}{S}

\title{%
  Lean\,4-Formalisierung der Durchlaufbedingung\\[0.3em]
  \large Verfügbare Mathlib-Bausteine, DC-relevante Lean-Snippets
  und offene Lücken im EABC-Rahmen%
}
\author{Thomas Hoffbauer}
\date{\today}

\begin{document}
\maketitle

\begin{abstract}
Die Durchlaufbedingung (DC) --- die kritische offene Lücke im oktonionischen
Collatz-Beweisversuch --- besagt, dass jede Collatz-Bahn nach höchstens~$m_0$
konsekutiven BC-Schritten zwingend einen EA-Schritt durchläuft.
Dieser Text untersucht, was davon in Lean\,4 / Mathlib bereits formalisierbar ist
und was nicht.
Abschnitt~\ref{sec:mathlib} inventarisiert die relevanten Mathlib-Werkzeuge
(2-adische Bewertung, Kongruenzrechnung, Syracuse-Abbildung).
Abschnitt~\ref{sec:projekt} beschreibt die Lean-Dateien des Projekts
und stellt den Stand realistisch dar.
Abschnitt~\ref{sec:snippets} enthält konkrete Lean\,4-Code-Snippets
für die drei DC-Bestandteile.
Abschnitt~\ref{sec:luecken} benennt ehrlich die verbleibenden Beweislücken.
\end{abstract}

\tableofcontents

\bigskip
\begin{tcolorbox}[warnung, title={Zur Einordnung}]
Die Collatz-Vermutung ist unbewiesen. Die DC ist ebenfalls unbewiesen und gilt
als mindestens ebenso schwer. Dieser Text unterscheidet konsequent zwischen
\emph{formal beweisbaren} modularen Fakten (die Lean kompiliert)
und \emph{offenen Lücken} (die \texttt{sorry} erfordern würden).
\end{tcolorbox}

% ============================================================
\section{Überblick: DC und EABC-Rahmen}
\label{sec:ueber}
% ============================================================

\subsection{Die Durchlaufbedingung}

\begin{tcolorbox}[kernidee]
\textbf{(DC) --- Durchlaufbedingung.}\quad
Für jedes ungerade $n \in \N$ existiert ein uniform beschränktes~$m_0$ und ein
$m \leq m_0$, so dass
\[
  \vii\!\bigl(3\,\U^m(n) + 1\bigr) \;\geq\; 2,
\]
wobei $\U(n) = (3n+1)/2^{\vii(3n+1)}$ die odd-to-odd-Abbildung bezeichnet.

\medskip
\textbf{Umformulierung:} Jede Collatz-Bahn durchläuft nach spätestens~$m_0$
konsekutiven BC-Schritten ($\vii = 1$) zwingend einen EA-Schritt ($\vii \geq 2$).
\end{tcolorbox}

\subsection{Die EABC-Klassifikation (mod~12)}

\begin{definition}[EABC-Typen]\label{def:eabc}
Für ungerades $n \in \N$:
\begin{align*}
  \text{E-Typ}:\quad n &\equiv 1 \pmod{12}, \quad \vii(3n+1) \geq 2,\\
  \text{A-Typ}:\quad n &\equiv 5 \pmod{12}, \quad \vii(3n+1) \geq 2,\\
  \text{B-Typ}:\quad n &\equiv 7 \pmod{12}, \quad \vii(3n+1) = 1,\\
  \text{C-Typ}:\quad n &\equiv 11 \pmod{12}, \quad \vii(3n+1) = 1.
\end{align*}
EA-Schritte sind Kontraktionsschritte, BC-Schritte sind Expansionsschritte.
\end{definition}

\subsection{Drei Teilaussagen der DC}

Das Dokument \texttt{collatz\_dc\_morley\_walter.tex} des Projekts formuliert
drei Teilaussagen:

\begin{enumerate}[label=(\roman*)]
  \item \textbf{Kein B~folgt auf B:}
    Wenn $n \equiv 7 \pmod{12}$ (B-Typ), dann ist
    $\U(n) \equiv 11 \pmod{12}$ (C-Typ), nie wieder B-Typ.
  \item \textbf{Exponentiell wachsende C-Ketten-Barriere:}
    C-Ketten der Länge~$k$ erfordern Startwerte in Residuenklassen
    der Dichte~$1/2^{k+2}$.
  \item \textbf{Geometrisch-heuristische Deutung:}
    Der Morley-Walter-Rahmen liefert Intuition, aber keinen Beweis der DC.
\end{enumerate}

Teilaussage~(i) ist formal beweisbar und gehört zum Abschnitt \ref{sec:b-nicht-b}.
Teilaussage~(ii) ist formal formulierbar, der Dichteabfall ist beweisbar,
die Schranke~$m_0$ hingegen nicht.
Teilaussage~(iii) ist rein heuristisch.

% ============================================================
\section{Mathlib-Werkzeuge: Was ist verfügbar?}
\label{sec:mathlib}
% ============================================================

\subsection{2-adische Bewertung}

\begin{tcolorbox}[ergebnis]
\textbf{In Mathlib verfügbar:}
\begin{itemize}
  \item \lstinline|padicValNat p n| --- p-adische Bewertung auf $\N$
    (\texttt{Mathlib.NumberTheory.Padics.PadicVal.Basic})
  \item \lstinline|padicValInt p n| --- p-adische Bewertung auf $\Z$
  \item \lstinline|Nat.trailingZeros n| --- Anzahl der abschließenden Nullbits
    in der Binärdarstellung; für $n > 0$ gilt
    \lstinline|Nat.trailingZeros n = padicValNat 2 n|
  \item Lemmata: \lstinline|padicValNat_dvd_iff_le|,
    \lstinline|pow_padicValNat_dvd|, \lstinline|padicValNat.self|
\end{itemize}
\end{tcolorbox}

Für die DC ist die zentrale Größe $\vii(3n+1) = \mathtt{padicValNat}\ 2\ (3n+1)$.

\subsection{Kongruenzrechnung}

\begin{tcolorbox}[ergebnis]
\textbf{In Mathlib verfügbar:}
\begin{itemize}
  \item \lstinline|ZMod n| --- Restklassen modulo~$n$ als Typ
    (\texttt{Mathlib.Data.ZMod.Basic})
  \item \lstinline|Int.ModEq n a b| --- Notation \lstinline|a == b [ZMOD n]|
    (\texttt{Mathlib.Data.Int.ModEq})
  \item \lstinline|omega| --- Taktik für lineare Arithmetik über $\Z$ und $\N$
  \item \lstinline|decide| --- Entscheidungsprozedur für endliche Kongruenzen
  \item Chinesischer Restsatz: \lstinline|ZMod.chineseRemainder|
\end{itemize}
\end{tcolorbox}

\subsection{Iteration und Periodische Bahnen}

\begin{tcolorbox}[ergebnis]
\textbf{Verfügbar (Lean\,4-Core + Mathlib):}
\begin{itemize}
  \item \lstinline|f^[n]| --- $n$-fache Iteration (\lstinline|Function.iterate|)
  \item \lstinline|Function.IsPeriodicPt f n x| ---
    $x$ ist periodischer Punkt der Periode~$n$
  \item \lstinline|Function.minimalPeriod| --- minimale Periode
\end{itemize}
\end{tcolorbox}

\begin{tcolorbox}[lücke]
\textbf{Nicht in Mathlib:}
\begin{itemize}
  \item Für die spezifische Collatz-/Syracuse-Abbildung gibt es in Mathlib
    keinen allgemeinen Satz über Zyklenfreiheit.
  \item Die Aussage
    ``$\forall\, n > 0,\ \forall\, k \geq 1,\ \Syr^{[k]}(n) \neq n$''
    (keine positiven periodischen Punkte) ist offen.
\end{itemize}
\end{tcolorbox}

\subsection{Externe Lean\,4-Projekte (Juni~2026)}

\begin{description}
  \item[\texttt{johnjanik/syracuse-confinement}]
    12\,947 Zeilen, 45 Dateien (Lean\,4 + Mathlib v4.27.0).
    Formale Reduktion der Collatz-Vermutung auf ein Diophantisches
    Confinement-Problem. Enthält \texttt{Syracuse.lean} mit
    \lstinline|syr_identity| (Tao-2022-Identität), \lstinline|val2|,
    \lstinline|syr_odd|, \lstinline|syracuse_descent_criterion|.
    \emph{Kritische Lücke:} \lstinline|nu3_linear_bound [sorry]|.
  \item[\texttt{ericmerle3789/collatz-conditional-cycles}]
    April~2026, 14\,110 Zeilen. Bedingte Nichtexistenz nicht-trivialer
    Collatz-Zyklen (Baker-Theorem als Axiom).
    Kein \lstinline|sorry| im Kernkode, aber externe arithmetische Bounds
    (\texttt{BakerSeparation}, \texttt{BarinaVerification}) als Axiome.
  \item[\texttt{shuanat/collatz-lean4}]
    Formalisierung von $\mathrm{ord}_{2^t}(3) = 2^{t-2}$,
    relevant für 2-adische Zyklenanalyse.
  \item[\texttt{QuixiAI/collatz}]
    Proof strategy framework: \lstinline|nextOdd_mod_lower/drop|,
    2-adische Instabilitäts-Lemmata; kritische Schritte als Axiome.
\end{description}

% ============================================================
\section{Vorhandene Lean-Bausteine im Projekt}
\label{sec:projekt}
% ============================================================

Das Projekt enthält folgende Lean-Dateien. Die ehrliche Einschätzung:

\begin{center}
\begin{tabular}{lll}
\toprule
Datei & Inhalt & DC-Relevanz \\
\midrule
\texttt{LeanTest1.lean} & \lstinline|and_commutative| (trivial) & keine \\
\texttt{LeanBeweis.lean} & Ramanujan-Petersson (alles \texttt{sorry}) & keine \\
\texttt{Betrtrand\_Lean/} & Leeres Hello-World-Projekt & keine \\
\texttt{ptolemaeus-lean/} & Ptolemäus-Ungleichung (ABC.lean, groß) & keine \\
\bottomrule
\end{tabular}
\end{center}

\begin{tcolorbox}[warnung]
Das Projekt enthält \emph{keinen} Lean-Code, der direkt auf Collatz
oder die DC anwendbar wäre.
Die mathematisch substanzielle Datei \texttt{Lean Vorschläge.tex}
behandelt 2HDM-Stabilität (nicht Collatz).
Die DC-relevanten informalen Beweise sind in
\texttt{collatz\_dc\_morley\_walter.tex} und
\texttt{collatz\_eabc\_spektrum.tex} dokumentiert.
\end{tcolorbox}

Die folgenden informalen Resultate aus den Projekt-LaTeX-Dateien
sind DC-relevant und prinzipiell formalisierbar:

\begin{enumerate}[label=(\alph*)]
  \item Proposition aus \texttt{collatz\_eabc\_spektrum.tex}:
    Für $n \equiv 7, 11 \pmod{12}$ gilt $\vii(3n+1) = 1$ (Ein-Halbierungs-Klassen).
  \item B-Typ~$\to$~C-Typ-Übergang (aus \texttt{collatz\_dc\_morley\_walter.tex}):
    Wenn $n \equiv 7 \pmod{12}$, dann $\U(n) \equiv 11 \pmod{12}$.
  \item C-Ketten-Dichte: Länge-$k$-C-Ketten erfordern Dichte~$1/2^{k+2}$.
\end{enumerate}

% ============================================================
\section{Lean\,4-Beweisvorschläge}
\label{sec:snippets}
% ============================================================

\subsection{Grunddefinitionen}

\begin{tcolorbox}[lean]
\begin{lstlisting}
import Mathlib.NumberTheory.Padics.PadicVal.Basic
import Mathlib.Data.ZMod.Basic
import Mathlib.Data.Int.ModEq

-- 2-adische Bewertung (Abkuerzung)
abbrev nu2 (n : Nat) : Nat := padicValNat 2 n

-- Syracuse-Abbildung (auf ungeraden Zahlen)
noncomputable def syr (n : Nat) : Nat :=
  (3 * n + 1) / 2 ^ nu2 (3 * n + 1)

-- EABC-Typen als Praedikate
def isEType (n : Nat) : Prop := n % 12 = 1
def isAType (n : Nat) : Prop := n % 12 = 5
def isBType (n : Nat) : Prop := n % 12 = 7
def isCType (n : Nat) : Prop := n % 12 = 11

-- BC-Schritt (Expansion): nu2 = 1
def isBCStep (n : Nat) : Prop := nu2 (3 * n + 1) = 1

-- EA-Schritt (Kontraktion): nu2 >= 2
def isEAStep (n : Nat) : Prop := nu2 (3 * n + 1) >= 2
\end{lstlisting}
\end{tcolorbox}

\subsection{B-Typ-Schritt und Halbierungstiefe~1}
\label{sec:b-nicht-b}

Dies ist die einfachste DC-relevante Aussage: B-Typ- und C-Typ-Zahlen
sind exakt die Ein-Halbierungs-Klassen. Die Beweise gelingen per
\lstinline|omega| oder \lstinline|decide|.

\begin{tcolorbox}[lean, title={Lean\,4-Code: BC-Schritte sind genau Halbierungstiefe~1}]
\begin{lstlisting}
-- Fuer n = 12k+7: 3n+1 = 2*(ungerade), also nu2 = 1
lemma bType_nu2_eq_one (k : Nat) :
    nu2 (3 * (12 * k + 7) + 1) = 1 := by
  simp [nu2, padicValNat]
  -- 3*(12k+7)+1 = 36k+22 = 2*(18k+11)
  -- 18k+11 ist ungerade (kongruent 1 mod 2), also nu2 = 1
  sorry -- lueckenfueller; omega schliesst die Arithmetik

-- Fuer n = 12k+11: analoges Argument
lemma cType_nu2_eq_one (k : Nat) :
    nu2 (3 * (12 * k + 11) + 1) = 1 := by
  sorry -- 3*(12k+11)+1 = 36k+34 = 2*(18k+17)
        -- 18k+17 ist ungerade, also nu2 = 1

-- Alternative: decide (fuer festes Modell)
example : forall r : ZMod 12,
    (r = 7 or r = 11) ->
    (3 * r + 1 : ZMod 12) = 10 := by
  decide
\end{lstlisting}
\end{tcolorbox}

\begin{bemerkung}
Die \lstinline|sorry|-Platzhalter markieren Stellen, an denen
\lstinline|padicValNat| direkt ausgerechnet werden muss.
Korrekt geht das so: Man zeigt zunächst
$2 \mid (3n+1)$ und $\neg\,(4 \mid (3n+1))$, woraus $\vii(3n+1) = 1$ folgt.
Mathlib stellt hierfür \lstinline|padicValNat.eq_one_of_dvd_not_pow_dvd| bereit.
\end{bemerkung}

\subsection{B-Typ kann nicht auf B-Typ folgen}

\begin{tcolorbox}[lean, title={Lean\,4-Code: B→C, kein B→B}]
\begin{lstlisting}
-- Der Kern: n = 12k+7 (B-Typ),
-- dann U(n) = (3n+1)/2 = 18k+11, also Typ C (kongruent 11 mod 12)
lemma b_maps_to_c (k : Nat) :
    syr (12 * k + 7) % 12 = 11 := by
  -- syr (12k+7) = (3*(12k+7)+1)/2 = (36k+22)/2 = 18k+11
  simp [syr, bType_nu2_eq_one]
  omega

-- Korollar: Keine reinen B-Ketten
corollary no_bb_chain (n : Nat) (hn : isBType n) :
    not (isBType (syr n)) := by
  obtain k := Nat.exists_eq_add_of_le (Nat.pos_of_ne_zero (by omega))
  simp [isBType, isCType] at *
  -- syr n kongruent 11 (mod 12), also kein B-Typ
  sorry

-- Staerker: Fuer C-Typ gilt syr(n) nicht kongruent 11 (mod 24)
-- (C->A oder C->E, nie wieder C auf mod-24-Ebene)
lemma c_maps_not_to_c_mod24 (k : Nat) :
    syr (24 * k + 11) % 24 != 11 := by
  sorry
\end{lstlisting}
\end{tcolorbox}

\begin{tcolorbox}[ergebnis]
Die Aussage ``B-Typ kann nicht auf B-Typ folgen'' ist
\textbf{formal beweisbar} mit \lstinline|omega| und
elementarer \lstinline|padicValNat|-Arithmetik.
Der Beweis verläuft in drei Schritten:
\begin{enumerate}
  \item Zeige $n \equiv 7 \pmod{12} \Rightarrow \vii(3n+1) = 1$
    (da $3n+1 \equiv 22 \equiv 2 \pmod{4}$).
  \item Berechne $\U(n) = (3n+1)/2 \equiv 11 \pmod{12}$ via \lstinline|omega|.
  \item Schließe: $\U(n)$ ist C-Typ, also kein B-Typ.
\end{enumerate}
Dies lässt sich in Lean\,4 vollständig ohne \lstinline|sorry| formalisieren.
\end{tcolorbox}

\subsection{C-Ketten-Dichte und exponentiell wachsende Barriere}

\begin{tcolorbox}[lean, title={Lean\,4-Code: Länge-k-C-Ketten und Restklassen-Dichte}]
\begin{lstlisting}
-- Definition: n bildet den Beginn einer C-Kette der Laenge k
-- wenn syr^[0](n), syr^[1](n), ..., syr^[k-1](n) alle C-Typ sind
def isCChainStart (n : Nat) (k : Nat) : Prop :=
  forall i, i < k -> isCType (syr^[i] n)

-- Fuer k=1: Aequivalent zu isCType n
lemma cchain_one (n : Nat) :
    isCChainStart n 1 <-> isCType n := by
  simp [isCChainStart, isCType]

-- Die Bedingung fuer eine C-Kette der Laenge 2:
-- Man braucht n kongruent 11 (mod 12) UND syr(n) kongruent 11 (mod 12)
-- syr(n) = (3n+1)/2: fuer n=12k+11 gilt syr(n)=18k+17 kongruent 5 (mod 12)
-- Also: syr(n) ist A-Typ! Kein C-Typ.
-- Daraus folgt: C-Ketten der Laenge >= 2 sind mod-12-unmoeglich.

-- Formaler Satz: Kein C folgt unmittelbar auf C (mod 12)
theorem no_c_chain_length_2_mod12 (n : Nat)
    (hn : isCType n) : not (isCType (syr n)) := by
  simp [isCType] at hn
  -- n kongruent 11 (mod 12) -> 3n+1 kongruent 10 (mod 12) -> (3n+1)/2 kongruent 5 (mod 12)
  -- 5 != 11, also syr(n) nicht kongruent 11 (mod 12)
  sorry -- Per padicValNat + omega; vollstaendig beweisbar

-- Auf mod-24-Ebene: Dichte 1/4 fuer Laenge-1-C-Ketten, 1/8 fuer Laenge-2
-- (verfeinert auf mod-48, mod-96, ...)
example : forall r : ZMod 12,
    r = (11 : ZMod 12) -> (3 * r + 1 : ZMod 12) = (10 : ZMod 12) := by
  decide

example : forall r : ZMod 12,
    r = (11 : ZMod 12) -> (3 * r + 1) / 2 = (5 : ZMod 12) := by
  decide
\end{lstlisting}
\end{tcolorbox}

\begin{bemerkung}
Die mod-12-Analyse zeigt etwas Stärkeres als erwartet:
Nicht nur ``B-Typ kann nicht auf B-Typ folgen'', sondern auch
\textbf{C-Typ kann nicht auf C-Typ folgen} (auf mod-12-Ebene).
Tatsächlich gilt: Wenn $n \equiv 11 \pmod{12}$ (C-Typ), dann
\[
  \U(n) = \frac{3n+1}{2} \;\equiv\; 5 \pmod{12} \quad (\text{A-Typ}).
\]
Damit sind C-Ketten auf mod-12-Ebene unmöglich!
Die ``Dichte~$1/2^{k+2}$''-Aussage des Projekts bezieht sich auf
eine feinere Analyse mod~$2^{k+2}$, die für spätere Elemente in einer Bahn gilt.
\end{bemerkung}

\subsection{Ansatz für die Endlichkeit von C-Ketten (2-adisch)}

Die stärkste verfügbare Strategie für die DC nutzt die 2-adische Instabilität
der Syracuse-Abbildung:

\begin{tcolorbox}[lean, title={Lean\,4-Code: 2-adische Instabilität (Skizze)}]
\begin{lstlisting}
-- Die 2-adische Erweiterung: Arbeiten in Z_2
-- (Mathlib: PadicInt 2, PadicIntegers)
-- Ziel: Zeige, dass S : Z_2 -> Z_2 keine positiven periodischen Punkte hat.

-- Zyklengleichung: Wenn syr^[k](n) = n fuer k Schritte, dann
-- n * (2^e - 3^k) = Summe von Potenzen
-- (e = Gesamtzahl der Halbierungsschritte)

-- Aus Mihailescu/Catalan: 2^e = 3^k + 1 nur fuer e=2, k=1 -> n=1 (trivial)
-- Alle anderen Loesungen: 2^e - 3^k < 0 -> negative ganze Zahlen

-- Formale Aussage (bedingt auf Baker-Theorem):
theorem no_positive_cycle
    (baker : forall e k : Nat, e >= 1 -> k >= 1 ->
             2^e != 3^k + 1 or (e = 2 and k = 1))
    (n : Nat) (hn : n > 1) (k : Nat) (hk : k >= 1) :
    syr^[k] n != n := by
  sorry -- Hauptluecke: Verbindet Zyklengleichung mit baker-Axiom

-- Die Tao-2022-Identitaet (aus johnjanik/syracuse-confinement):
-- syr^[k](n) * 2^|a^(k)| = 3^k * n + G_k
-- wobei G_k = Summe 3^(k-i) * 2^|a^(i)|
-- Dies ist das Herzstueck jeder Zyklenanalyse.
theorem tao_identity (n k : Nat) (hn : Odd n) :
    exists Gk : Nat,
      syr^[k] n * 2^k = 3^k * n + Gk := by
  sorry -- In johnjanik/syracuse-confinement formalisiert (syr_identity)
\end{lstlisting}
\end{tcolorbox}

% ============================================================
\section{Offene Lücken: Was ist in Lean nicht beweisbar?}
\label{sec:luecken}
% ============================================================

\subsection{Ehrliche Bestandsaufnahme}

\begin{tcolorbox}[lücke]
\textbf{Was in Lean\,4 / Mathlib (Stand Juni~2026) \emph{nicht} ohne \texttt{sorry} beweisbar ist:}

\begin{enumerate}[label=(\arabic*)]
  \item \textbf{Die vollständige DC:}
    Der Satz $\forall n \in \N,\ \exists m_0,\ \exists m \leq m_0:\
    \vii(3\,\U^m(n)+1) \geq 2$
    ist unbewiesen und würde die Collatz-Vermutung implizieren.

  \item \textbf{Die gleichmäßige Schranke $m_0$:}
    Selbst wenn C-Ketten der Länge~$k$ exponentiell selten sind,
    folgt daraus keine global uniforme Schranke.

  \item \textbf{Zyklenfreiheit positiver Bahnen:}
    ``Alle positiven periodischen Punkte der Syracuse-Abbildung
    sind der Fixpunkt~$1$.''
    Dies ist bedingt formalisiert (Baker-Theorem als Axiom in
    \texttt{ericmerle3789/collatz-conditional-cycles}),
    aber nicht unconditional beweisbar.

  \item \textbf{Positive 2-adische Bahnen:}
    Die Aussage, dass jede positive ganzzahlige Bahn
    unter $\Syr$ beschränkt bleibt und~$1$ erreicht,
    ist äquivalent zur Collatz-Vermutung.

  \item \textbf{Morley-Walter-Ansatz:}
    Die geometrische Interpretation via Morley-Walter-Rahmen
    liefert Heuristik, aber kein formal verifizierbares Argument.
\end{enumerate}
\end{tcolorbox}

\subsection{Was noch in Mathlib fehlt}

Für eine vollständige Formalisierung der DC müsste in Mathlib noch
folgendes entwickelt werden:

\begin{enumerate}[label=(\alph*)]
  \item \textbf{2-adische Zahlen $\Ztwo$ als dynamischer Raum:}
    Zwar ist \lstinline|PadicInt 2| in Mathlib vorhanden,
    aber die ergodische Theorie der Collatz-Abbildung auf $\Ztwo$
    (nach Lagarias/Terras) ist nicht formalisiert.

  \item \textbf{Uniform Distribution modulo Zweierpotenzen:}
    Taos 2022-Ergebnis (``fast alle Bahnen fallen unter $n^\epsilon$'')
    ist in Lean nirgends formalisiert.

  \item \textbf{Baker-Schranken für $|2^e - 3^k|$:}
    Die linear independence of logarithms (Baker 1966) wäre
    als Mathlib-Axiom nötig --- derzeit nur als
    \lstinline|BakerSeparation|-Axiom in externen Projekten.

  \item \textbf{Ergodische Theorie der Collatz-Abbildung:}
    Die Messinvarianz der 2-adischen Haar-Maßes unter~$\Syr$
    und die Ergodizität sind in Mathlib nicht vorhanden.
\end{enumerate}

\subsection{Realistischer Status der Lean-Formalisierbarkeit}

\begin{center}
\begin{tabular}{p{5.5cm}cc}
\toprule
Aussage & Lean-Status & Referenz \\
\midrule
B-Typ $\Rightarrow$ C-Typ (kein B$\to$B) & \textcolor{bewiesengrün}{\textbf{beweisbar}} & \S\ref{sec:b-nicht-b} \\
C-Typ $\Rightarrow$ A-Typ (kein C$\to$C) & \textcolor{bewiesengrün}{\textbf{beweisbar}} & \S\ref{sec:snippets} \\
$\vii(3n+1)=1$ für $n \equiv 7,11 \pmod{12}$ & \textcolor{bewiesengrün}{\textbf{beweisbar}} & \S\ref{sec:snippets} \\
$\vii(3n+1)\geq 2$ für $n \equiv 1,5 \pmod{12}$ & \textcolor{bewiesengrün}{\textbf{beweisbar}} & \S\ref{sec:snippets} \\
Tao-Identität $\Syr^{[k]}(n)\cdot 2^e = 3^k n + G_k$ & \textcolor{warnorange}{\textbf{extern}} & johnjanik \\
Keine positiven Zyklen (bedingt) & \textcolor{warnorange}{\textbf{bedingt}} & ericmerle, baker \\
Keine positiven Zyklen (unbedingt) & \textcolor{lückenrot}{\textbf{offen}} & --- \\
DC: Gleichmäßige Schranke $m_0$ & \textcolor{lückenrot}{\textbf{offen}} & --- \\
Vollständige Collatz-Vermutung & \textcolor{lückenrot}{\textbf{offen}} & --- \\
\bottomrule
\end{tabular}
\end{center}

% ============================================================
\section{Fazit und Empfehlungen}
\label{sec:fazit}
% ============================================================

\subsection{Was im Projekt formalisiert werden sollte}

Wenn das Projekt eine Lean-Formalisierung anstrebt, ist folgende
Priorisierung realistisch:

\begin{enumerate}[label=\textbf{Priorität~\arabic*:}]
  \item \textbf{Elementare EABC-Lemmas:}
    Die mod-12-Eigenschaften (B$\to$C, C$\to$A, E$\to$E-oder-A,
    A$\to$E-oder-B) sind \emph{vollständig beweisbar} und bilden
    den formalen Kern der EABC-Klassifikation.
    Zeitaufwand: einige Stunden.

  \item \textbf{Ein-Halbierungs-Satz:}
    $\vii(3n+1) = 1 \Leftrightarrow n \equiv 7$ oder $n \equiv 11 \pmod{12}$.
    Beweisbar mit \lstinline|padicValNat| + \lstinline|omega|.
    Zeitaufwand: ein Tag.

  \item \textbf{Tao-Identität:}
    Anstatt das Rad neu zu erfinden, sollte die bestehende
    Formalisierung von \texttt{johnjanik/syracuse-confinement}
    importiert oder als Referenz verwendet werden.

  \item \textbf{Ehrlichkeit über die DC:}
    Die DC ist in Lean nicht beweisbar ohne einen Durchbruch,
    der entweder der Collatz-Vermutung entspricht oder sie impliziert.
    Jede Lean-Datei, die die DC ``beweist'', enthält notwendigerweise
    entweder \lstinline|sorry| oder falsche Axiome.
\end{enumerate}

\subsection{Abschließende Einschätzung}

\begin{tcolorbox}[kernidee]
\textbf{Was Lean\,4 leisten kann:}
\begin{itemize}
  \item Scharfe formale Dokumentation der EABC-Struktur
    und der mod-12-Barrieren (vollständig formalisierbar).
  \item Präzise Benennung der Lücken: Welche Axiome würden die DC implizieren?
  \item Verifikation der elementaren Schritttypen und Übergangsregeln.
\end{itemize}

\textbf{Was Lean\,4 nicht kann:}
\begin{itemize}
  \item Den Beweis der DC oder der Collatz-Vermutung schließen.
  \item Die heuristische Morley-Walter-Intuition in Theoreme übersetzen,
    ohne tiefe Hilfsmittel (Ergodentheorie, Baker-Schranken).
\end{itemize}

Die DC bleibt die tiefste offene Lücke. Ihre formale Schließung
würde einen mathematischen Durchbruch darstellen.
\end{tcolorbox}

\end{document}
